\documentclass[conference]{IEEEtran}
\IEEEoverridecommandlockouts
% The preceding line is only needed to identify funding in the first footnote. If that is unneeded, please comment it out.
\usepackage{cite}
\usepackage{amsmath,amssymb,amsfonts}
\usepackage{algorithmic}
\usepackage{graphicx}
\usepackage{textcomp}
\usepackage{xcolor}
\usepackage{booktabs}
\def\BibTeX{{\rm B\kern-.05em{\sc i\kern-.025em b}\kern-.08em
    T\kern-.1667em\lower.7ex\hbox{E}\kern-.125emX}}
\begin{document}

\title{Learnable Pair-to-Edge Gating of Frozen 3D Pair Representations\\
for Spectral Molecular Property Prediction}

% TODO(authors): 作者名单待与学姐对齐后填写真实姓名/单位。
% 现在保持匿名占位 —— ICDE 研究论文投稿是双盲评审,匿名是安全默认。
\author{\IEEEauthorblockN{Anonymous Author(s)}
\IEEEauthorblockA{\textit{Paper submitted for double-blind review}}}
% --- 确定作者后,用下面的模板替换上面的匿名块: ---
% \author{
% \IEEEauthorblockN{1\textsuperscript{st} Given Name Surname}
% \IEEEauthorblockA{\textit{dept. name of organization (of Aff.)} \\
% \textit{name of organization (of Aff.)}\\
% City, Country \\
% email address or ORCID}
% \and
% \IEEEauthorblockN{2\textsuperscript{nd} Given Name Surname}
% \IEEEauthorblockA{\textit{dept. name of organization (of Aff.)} \\
% \textit{name of organization (of Aff.)}\\
% City, Country \\
% email address or ORCID}
% }

\maketitle

\begin{abstract}
Two-dimensional graph neural networks for molecular property prediction capture bond topology well but cannot directly access through-space geometric interactions, while purely three-dimensional models pay a steep cost in conformer dependence and compute. We study a lightweight middle ground: starting from a Chebyshev~II spectral graph network as the two-dimensional backbone, we inject a frozen Uni-Mol pair representation as a single scalar gate on each chemical bond, with a bias initialization $b{=}{+}5$ that places the initial gate at $\sigma(5)\approx 0.993$ and preserves the unweighted Laplacian at epoch zero. The resulting variant, \emph{V2-T5}, ties the two-dimensional baseline V0 within seed variance at $n{=}9$ training seeds on BBBP ($0.831 \pm 0.032$ vs.\ $0.828 \pm 0.027$); a more expressive dynamic variant, \emph{T6}, that refines pair and node features at each layer ties it as well ($0.831$), so under our small-data contrastive-pretraining regime the extra capacity buys no test-time gain over the static near-identity gate. A post-hoc audit of the trained gates shows that on BBBP V2-T5 learns a chemistry-aware routing that up-weights aromatic bonds and down-weights saturated single bonds (Pearson $+0.66$ across seeds), whereas on BACE the same architecture learns no detectable per-bond-type preference. On FreeSolv, V2-T5 reduces test RMSE from $0.760$ to $0.675$, but using a Gaussian-basis distance surrogate in place of the Uni-Mol pair tensor (\S\ref{sec:v2t5}), so this is an indirect test only. We additionally report a fingerprint-only control and a 30-seed random forest baseline that exceed the two-dimensional graph baseline on BACE by $8.9$ and $13.7$ AUC points respectively, and argue that such controls should be standard practice when evaluating graph-based molecular encoders.
\end{abstract}

\begin{IEEEkeywords}
molecular property prediction, graph neural networks, spectral methods, 3D molecular representation, self-supervised learning
\end{IEEEkeywords}

% ---------------------------------------------------------------------
\section{Introduction}
\label{sec:intro}

Molecular property prediction is a central problem in drug discovery and computational chemistry, and a long line of work has shown that accurate predictions hinge on representations that capture both the topology of chemical bonds and the geometry of three-dimensional conformers. Two-dimensional graph neural networks exploit bond-level structure but cannot directly access through-space interactions such as steric clashes, hydrogen bonds, or favorable stacking. Conversely, purely three-dimensional models depend on the quality of conformer generation and tend to be computationally expensive at the scale of high-throughput screening. A practical alternative is to start from a strong two-dimensional backbone and inject three-dimensional information only where it is most useful, but the design space for such an injection remains underexplored.

Self-supervised pretraining has produced increasingly capable molecular encoders along both modalities. On the two-dimensional side, contrastive frameworks such as MolCLR~\cite{wang2022molclr} and graph transformers such as GROVER~\cite{rong2020grover} learn transferable atom and bond representations from millions of unlabeled molecules. On the three-dimensional side, Uni-Mol~\cite{zhou2023unimol} pretrains an SE(3)-equivariant encoder over conformer ensembles and produces, in addition to atomic embeddings, a dense pair representation $P \in \mathbb{R}^{N \times N \times d_p}$ that encodes geometric relationships between every atom pair. Multi-view pretraining approaches such as GraphMVP~\cite{liu2022graphmvp} have further shown that aligning two- and three-dimensional views during pretraining benefits downstream prediction. The complementary question of how to consume a frozen three-dimensional pair representation at downstream time, without retraining the upstream encoder, has received considerably less attention.

Spectral graph neural networks offer a particularly clean entry point for such an injection. The recently introduced ChebNet II~\cite{he2022chebnetii} parametrizes graph convolutions as Chebyshev polynomials of the symmetric normalized Laplacian and admits per-edge weights as a first-class input. This makes the edge weights of the Laplacian a natural locus at which to fold in three-dimensional information: an off-the-shelf pair representation can be projected to a scalar gate on each chemical bond, and the resulting weighted Laplacian then drives the spectral filtering. The challenge is that naively initializing such a gate to small or random values immediately perturbs the spectrum away from the unweighted baseline, destabilizing the contrastive pretraining objective and erasing the inductive bias of the two-dimensional backbone before any useful learning has taken place.

We address this with two designs that share the same backbone---a dual-path Chebyshev II encoder with low-pass and high-pass towers, fused with a fingerprint MLP and pretrained by an NT-Xent contrastive loss across four views---but differ in how three-dimensional pair information is fed into the spectral filter. The first, \textbf{V2-T5}, applies a small MLP followed by a symmetrized sigmoid with bias initialization $b{=}{+}5$ to the Uni-Mol pair representation, yielding an initial per-bond gate of approximately $0.993$. This near-identity initialization preserves the two-dimensional baseline at epoch zero and lets the gate gradually move away from one only when the gradient supports it. The second, \textbf{T6}, replaces this static gate with a bidirectional pair-node update that iteratively refines both representations at each layer, trading stability for additional capacity. We additionally report a fingerprint-only control that uses no graph encoder, because such a control is frequently omitted from molecular benchmarks but turns out to be highly competitive.

To compare the three variants fairly we adopt the Uni-Mol scaffold fold split on the classification benchmarks and a scaffold split on FreeSolv, with matched training budgets and reporting format across all runs (BBBP at $n{=}9$ training seeds $\{9, 19, 29, 39, 49, 59, 69, 79, 89\}$; BACE and FreeSolv at $n{=}3$ seeds $\{9, 19, 29\}$). Under this protocol the static V2-T5 gating attains a mean test AUC of $0.831 \pm 0.032$ on BBBP, tied with the V0 baseline at $0.828 \pm 0.027$ and with T6 at $0.831 \pm 0.022$; on FreeSolv V2-T5 reduces test RMSE from $0.760$ to $0.675$ with a Gaussian-basis pair-feature surrogate rather than a Uni-Mol pair tensor, the strongest matched-seed test-time effect we observe. On BACE the fingerprint-only control reaches $0.846$ AUC and surpasses the two-dimensional graph baseline at $0.757$ by nearly nine AUC points, underscoring that strong fingerprint baselines should not be omitted from molecular benchmarks. The mechanism finding, by contrast, is robust: across every audited V2-T5 seed on BBBP the trained gate up-weights aromatic bonds and down-weights saturated single bonds (Pearson with the aromatic-bond indicator in $[0.6, 0.7]$), whereas on BACE the same architecture learns no detectable per-bond-type preference. Taken together, these results suggest that a frozen three-dimensional pair representation, when injected near identity and allowed to remain close to it, learns chemistry-aware bond routing on benchmarks where the relevant property has a chemistry signal not already captured by a fingerprint.

\noindent\textbf{Contributions.}
\begin{itemize}
  \item A static pair-to-edge gating mechanism (V2-T5) that projects a frozen Uni-Mol pair representation onto the chemical-bond edge weights of a Chebyshev II spectral GNN, using a near-identity initialization that preserves contrastive pretraining stability.
  \item A dynamic bidirectional pair-node update variant (T6) evaluated under matched compute, together with a per-seed analysis of where its additional optimization variance shows up under a three-seed BBBP protocol.
  \item A strict fair-comparison protocol on the Uni-Mol scaffold fold split with matched seeds, training budgets, and reporting format, including a fingerprint-only control that we find essential for honest assessment of graph-based molecular encoders.
\end{itemize}

% ---------------------------------------------------------------------
\section{Related Work}
\label{sec:related}

\subsection{Spectral and Message-Passing GNNs}
Spectral GNNs trace back to the Chebyshev formulation of Defferrard et al.~\cite{defferrard2016chebnet}, simplified by Kipf and Welling~\cite{kipf2017gcn}, and complemented by attention-based \texttt{GAT}~\cite{velickovic2018gat} and injective \texttt{GIN}~\cite{xu2019gin}. A complementary line revisits which polynomial basis yields the most expressive filter: Bernstein with non-negativity~\cite{he2021bernnet}, generalized PageRank~\cite{chien2021gprgnn}, and Chebyshev~II~\cite{he2022chebnetii}, the basis we adopt. An early molecular use of this adaptive-edge-spectral idea is the consistent-edge-attention spectral GCN of Shang et al.~\cite{shang2021eagcn}, which lets attention modulate the Laplacian; we follow the same move but route the per-edge weight through a Uni-Mol pair representation. Recently, Chen et al.~\cite{chen2024polygcl} extended polynomial filters to graph contrastive learning, and Nogueira et al.~\cite{nogueira2025spectra} used the eigenspectrum for target-aware augmentation. For molecular property prediction itself, Gilmer et al.~\cite{gilmer2017mpnn} unified earlier models under message-passing, and \texttt{AttentiveFP}~\cite{xiong2020attentivefp} remains a frequent attention-aggregation baseline; both operate on bond topology only, motivating the through-space gating studied here.

\subsection{3D Representations and Self-Supervised Pretraining}
Three-dimensional models inject coordinate information via continuous-filter convolutions~\cite{schutt2017schnet} (also the conceptual ancestor of our GBF FreeSolv-side control), directional angles~\cite{klicpera2020dimenet}, or $E(n)$ equivariance~\cite{satorras2021egnn}; \texttt{Uni-Mol}~\cite{zhou2023unimol} pretrains an SE(3)-equivariant encoder and exposes a dense pair tensor $P \in \mathbb{R}^{N\times N\times d_p}$ that we consume frozen, and \texttt{Frad}~\cite{feng2024frad} extends 3D pretraining with anisotropic-noise denoising. On the self-supervised side, \texttt{MolCLR}~\cite{wang2022molclr}, \texttt{GROVER}~\cite{rong2020grover}, and \texttt{GraphMVP}~\cite{liu2022graphmvp} cover 2D contrastive, graph-transformer, and 2D--3D alignment; \texttt{GEM}~\cite{fang2022gem}, \texttt{3D Infomax}~\cite{stark2022infomax}, \texttt{Mole-BERT}~\cite{xia2023molebert}, and \texttt{KPGT}~\cite{li2023kpgt} progressively add geometry-aware objectives, mutual-information maximization, and RDKit-knowledge fusion. More recent multi-task pretraining at moderate scale (\texttt{MiniMol}~\cite{klaeser2024minimol}) rivals foundation models on TDC ADMET; on the cheminformatics side, \texttt{MapLight}~\cite{notwell2023maplight} and \texttt{NovoExpert}-2~\cite{novoexpert2026} show that fingerprint + gradient-boosted trees top the TDC ADMET leaderboard, reinforcing that a strong fingerprint baseline is harder to beat than the graph-side literature reports. Our work is complementary: we do not propose a new SSL objective but a downstream-time mechanism for plugging a frozen 3D pair representation into a 2D spectral backbone, evaluated under the matched-protocol regime advocated by the \texttt{MoleculeNet} benchmark~\cite{wu2018moleculenet}.

% ---------------------------------------------------------------------
\section{Method}
\label{sec:method}

\subsection{Notation}
Let $G=(\mathcal{V},\mathcal{E})$ be a molecular graph with $N=|\mathcal{V}|$ heavy atoms and chemical bonds $\mathcal{E}\subseteq\mathcal{V}\times\mathcal{V}$. Each atom carries a feature vector $\mathbf{x}_i\in\mathbb{R}^{d_x}$ (atomic number, degree, hybridization, etc.); each bond an edge attribute $\mathbf{e}_{ij}\in\mathbb{R}^{d_e}$. We additionally consume a frozen pair representation $\mathbf{P}\in\mathbb{R}^{N\times N\times d_p}$ produced by a pretrained Uni-Mol encoder~\cite{zhou2023unimol}, indexed over all atom pairs $(i,j)$. The chemical-bond adjacency $A\in\{0,1\}^{N\times N}$ is symmetric (each bond contributes both directed edges $(i,j)$ and $(j,i)$) and induces the symmetric normalized Laplacian $L_{\mathrm{sym}}=I-D^{-1/2}AD^{-1/2}$, where $D$ is the degree matrix.

\subsection{Backbone: Dual-Path Chebyshev Spectral GNN}
Our backbone, denoted \texttt{LH\_Direct}, processes $G$ along two parallel spectral paths and fuses them with a molecular fingerprint stream. Each spectral path is a Chebyshev~II encoder~\cite{he2022chebnetii} that applies an order-$K$ polynomial filter to the normalized Laplacian:
\begin{equation}
    h^{(\ell)} = \sum_{k=0}^{K} \theta_k\, T_k(\tilde L)\, h^{(\ell-1)},
    \qquad \tilde L = 2L_{\mathrm{sym}}/\lambda_{\max} - I,
    \label{eq:chebnetii}
\end{equation}
where $T_k$ is the $k$-th Chebyshev polynomial and the coefficients $\{\theta_k\}_{k=0}^{K}$ are learnable. We fix $\lambda_{\max}{=}2$ for $\tilde L$ to avoid per-batch eigenvalue estimation. The two paths use independent coefficients tuned, by construction, towards low-pass and high-pass behavior respectively. The pooled graph representations $h_l, h_h\in\mathbb{R}^{d_h}$ are concatenated with a fingerprint embedding $h_{\mathrm{fp}}=\mathrm{MLP}(\mathrm{FP}(G))$, where $\mathrm{FP}$ concatenates PubChem (881), MACCS (167), and ErG (441) fingerprints into a $1489$-dimensional vector. The fused representation $z = [h_l; h_h; h_{\mathrm{fp}}]$ is the only object exposed to the downstream head.

\subsection{Variant V0: Two-Dimensional Baseline}
The V0 baseline runs \texttt{LH\_Direct} with $A$ unchanged, i.e.\ all bonds carry equal Laplacian weight. V0 ignores the Uni-Mol pair representation entirely and serves as the strict two-dimensional reference.

\subsection{Variant V2-T5: Static Pair-to-Edge Gating}
\label{sec:v2t5}
V2-T5 augments V0 with a learnable scalar gate on each chemical bond, computed from the frozen Uni-Mol pair representation. Concretely, for each bond $(i,j)\in\mathcal{E}$ we form
\begin{equation}
    g_{ij} = \sigma\!\left( \tfrac{1}{2}\bigl(\mathrm{MLP}_\phi(\mathbf{P}_{ij}) + \mathrm{MLP}_\phi(\mathbf{P}_{ji})\bigr) + b \right),
    \label{eq:v2t5gate}
\end{equation}
where $\mathrm{MLP}_\phi$ is a two-layer scalar-output network shared across edges, the symmetrization ensures $g_{ij}=g_{ji}$, and $\sigma$ is the logistic sigmoid. Crucially we initialize the bias term $b=+5$, which gives an initial gate of $\sigma(5)\approx 0.993$ across all bonds. This near-identity initialization preserves the V0 dynamics at epoch zero, so the contrastive pretraining objective starts from a near-optimal point and is free to move $g_{ij}$ away from $1$ only when gradient signal justifies it. The reweighted Laplacian $L_{\mathrm{sym}}(g)$ is then plugged into \eqref{eq:chebnetii} in place of $L_{\mathrm{sym}}$.

\noindent\textbf{Pair-representation source.} On the classification benchmarks (BBBP, BACE) we feed the frozen Uni-Mol pair representation directly into $\mathrm{MLP}_\phi$ in \eqref{eq:v2t5gate}. In the absence of a Uni-Mol scaffold-fold pipeline tailored to regression splits at the time of writing, the FreeSolv runs use an RDKit-generated 3D conformer with a Gaussian-basis-function (GBF) expansion of pairwise distances (64-dim, max distance $10$~\AA, $\sigma{=}0.5$) as a pair-feature surrogate; the gating mechanism, MLP head, bias initialization, and contrastive objective are otherwise identical. We are explicit about this asymmetry because the FreeSolv result is therefore an indirect rather than a direct test of the Uni-Mol contribution and should be read as such.

\noindent\textbf{Beyond a scalar gate.} The scalar gate of \eqref{eq:v2t5gate} collapses the $d_p$-dimensional pair representation onto a single number per bond, which is the smallest perturbation of the V0 Laplacian compatible with consuming any pair information at all. A natural generalization would replace the scalar gate with a low-rank tensor gate that emits a separate scalar for the low-pass and high-pass Chebyshev towers; we leave this to follow-up work (\S\ref{sec:conclusion}).

\subsection{Variant T6: Dynamic Pair-Node Update}
\label{sec:t6}
T6 trades the static gate of V2-T5 for a bidirectional pair-node update applied at each Chebyshev layer. Given current node features $h^{(\ell)}$ and the current pair bias $\mathbf{P}^{(\ell)}$, a multi-head attention block (four heads, two stacked layers, residual scaling $\alpha{=}0.1$) proceeds in two coupled steps. First, it computes query, key, and value projections $Q,K,V$ from $h^{(\ell)}$ and uses the pair bias as an additive term inside the softmax attention,
\begin{equation*}
    \mathrm{Attn}(Q,K,V;\mathbf{P}^{(\ell)}) = \mathrm{softmax}\!\bigl(QK^\top/\sqrt{d_h} + \mathbf{P}^{(\ell)}\bigr)\,V,
\end{equation*}
producing an updated node representation $h^{(\ell+1)}$ via the standard residual-and-LayerNorm wrapper. Second, the pre-softmax logits $QK^\top/\sqrt{d_h}$ are passed through a head-channel mixing projection to produce a pair residual $\Delta\mathbf{P}^{(\ell)}$ that is added back into the pair bias with the small residual scale $\alpha$,
\begin{equation*}
    \mathbf{P}^{(\ell+1)} = \mathrm{LayerNorm}\!\left(\mathbf{P}^{(\ell)} + \alpha\,\Delta\mathbf{P}^{(\ell)}\right).
\end{equation*}
The updated pair bias is then projected to edge gates via the same \eqref{eq:v2t5gate} machinery used by V2-T5. The small residual scaling $\alpha{=}0.1$ places T6 close to the identity at initialization---so that T6 nearly reduces to V2-T5---but unlike the static-gate special case, T6 has the additional capacity to update both representations as training proceeds. The price is one extra attention matmul per layer and additional optimization variance, which we examine empirically in Section~\ref{sec:experiments}.

\subsection{Training Objective and Downstream Evaluation}
\label{sec:training}
\noindent\textbf{Self-supervised pretraining.} All variants are pretrained with a four-view NT-Xent contrastive objective that pulls together low-pass, high-pass, mixed, and fingerprint embeddings of the same molecule and pushes apart embeddings of different molecules within a batch. We use temperature $\tau{=}0.5$, batch size $512$, and early-stopping based on training loss with patience $100$ epochs.

\noindent\textbf{Downstream linear probe.} After pretraining, the encoder is frozen and a two-layer logistic regression head (with a single Dropout layer set to $p{=}0.2$ during training and disabled during evaluation) is trained on top of the fused representation $z$. For each pretraining seed we run three independent evaluation restarts with probe-init seeds $\{9,19,29\}$ and report the test metric at the validation-best epoch. To prevent early-epoch noise from locking model selection on a lucky-spike checkpoint, we require both a minimum of $200$ probe epochs before any selection update and a strict improvement tolerance of $10^{-4}$.

% ---------------------------------------------------------------------
\section{Experiments}
\label{sec:experiments}

\subsection{Setup}
\noindent\textbf{Datasets.} BBBP (binary classification, 2{,}039 molecules), BACE (binary classification, 1{,}513 molecules), FreeSolv (regression, 642 molecules). All splits follow the Uni-Mol scaffold fold protocol.

\noindent\textbf{Seeds and budget.} BBBP is reported over nine training seeds $\{9,19,29,39,49,59,69,79,89\}$ and BACE and FreeSolv over three training seeds $\{9,19,29\}$ for the contrastive pretraining stage; the BBBP extension to $n{=}9$ is a dedicated replication discussed in \S\ref{sec:experiments}. For each pretrained checkpoint we run three independent downstream evaluation restarts and average over them to obtain one score per training seed. Pretraining runs for up to $1000$ epochs (clipped to $500$ when launched under the four-hour HPC QOS wall-clock cap) with batch size $512$ and NT-Xent temperature $\tau{=}0.5$, and is early-stopped on training loss with patience $100$. The downstream linear probe runs for at most $2000$ epochs with the validation-best checkpoint selected after the warm-up period of $200$ epochs described in Section~\ref{sec:training}. All variants share this protocol so that any difference in test metric is attributable to the encoder rather than to a difference in training budget or selection rule.

\noindent\textbf{Baselines.} V0 (2D-only), FP-only (fingerprint MLP head only, no graph encoder).

\noindent\textbf{Compute and hyperparameters.} Each pretraining run uses a single NVIDIA RTX 5000 Ada (32~GB) on the MBZUAI HPC cluster with approximately four hours of wall-clock per (variant, seed); the downstream linear probe runs in well under an hour. No hyperparameters are tuned per variant: all variants inherit $K{=}10$, hidden dimension $d_h{=}512$, dropout $0.2$, and the optimizer and scheduler of the V0 baseline. Code, processed datasets, and the corrected data rebuild scripts will be released alongside the camera-ready version.

\subsection{Main Results}

\begin{table*}[t]
\caption{Main results across molecular property prediction benchmarks. Classification reports test AUC (higher is better); regression reports test RMSE (lower is better). Classification benchmarks use the Uni-Mol scaffold-fold split; FreeSolv uses a standard scaffold split. Deep-model cells (V0, FP-only, V2-T5, T6) are mean~$\pm$~sample standard deviation across training seeds: $n{=}9$ for BBBP and $n{=}3$ for BACE and FreeSolv (see \S\ref{sec:experiments}). For variants with multiple downstream restarts per training seed, the per-seed mean (averaging over restarts) is the unit used in the std calculation, so inter-seed and intra-seed noise are not conflated. The RF row uses the 30-seed protocol of Deng et al.~\cite{deng2023systematic} on the matched Uni-Mol fold split ($n{=}30$). Dashes mark configurations that were not run (FP-only on BBBP and FreeSolv; T6 on FreeSolv).}
\label{tab:main}
\centering
\begin{tabular}{lccc}
\toprule
 & \multicolumn{2}{c}{Classification (AUC $\uparrow$)} & Regression (RMSE $\downarrow$) \\
\cmidrule(lr){2-3} \cmidrule(lr){4-4}
Variant & BBBP & BACE & FreeSolv \\
\midrule
V0 & 0.828 $\pm$ 0.027 & 0.757 $\pm$ 0.047 & 0.760 $\pm$ 0.058 \\
FP-only & -- & 0.846 $\pm$ 0.012 & -- \\
RF & 0.799 $\pm$ 0.007 & 0.894 $\pm$ 0.003 & 0.675 $\pm$ 0.010 \\
Chemprop & 0.815 $\pm$ 0.014 & 0.840 $\pm$ 0.045 & 0.879 $\pm$ 0.085 \\
V2-T5 & 0.831 $\pm$ 0.032 & 0.763 $\pm$ 0.010 & 0.675 $\pm$ 0.040 \\
T6 & 0.831 $\pm$ 0.022 & 0.768 $\pm$ 0.088 & -- \\
\bottomrule
\end{tabular}
\end{table*}

Table~\ref{tab:main} summarizes performance across the three benchmarks. On BBBP we now report $n{=}9$ training seeds for all deep variants ($\{9, 19, 29, 39, 49, 59, 69, 79, 89\}$). The static pair-to-edge gate (V2-T5) reaches $0.831 \pm 0.032$ mean test AUC, statistically indistinguishable from the two-dimensional baseline V0 ($0.828 \pm 0.027$) and from the dynamic pair-node update T6 ($0.831 \pm 0.022$); the three deep variants tie within seed variance on this benchmark. The original $n{=}3$ subsample (seeds $9, 19, 29$) placed V2-T5 at $0.862$ and V0 at $0.836$, a $+2.6$ AUC point gap that we now read as upper-tail sampling on the original seed set (\S\ref{sec:discussion}); we therefore foreground the mechanism finding (chemistry-aware aromatic-routing of the trained gate, \S\ref{sec:experiments}) rather than the headline AUC number. On FreeSolv, V2-T5 with the GBF pair-feature surrogate described in \S\ref{sec:v2t5} reduces test RMSE from $0.760$ to $0.675$ at $n{=}3$, an $11{.}2\%$ relative reduction; V2-T5 wins on all three matched seeds (mean delta $0.086$ RMSE, sign-test $p{=}0.25$), the strongest matched-seed evidence in the paper. While this number is not directly comparable to the Uni-Mol-driven BBBP result, it offers convergent evidence that 3D pair features---even a relatively coarse RDKit-distance encoding---contribute geometry information that the two-dimensional Laplacian cannot recover from topology alone. On BACE the fingerprint-only control attains $0.846 \pm 0.012$ AUC, exceeding V0's $0.757 \pm 0.047$ by $8.9$ points with substantially lower seed variance; this is the most informative single number in the table because it shows that on this benchmark a strong cheminformatics fingerprint, with no graph encoder at all, is harder to beat than the graph-side molecular literature usually assumes. The BACE V2-T5 and T6 entries are: V2-T5 reaches $0.763 \pm 0.010$ (with the tightest seed variance of any graph variant on the benchmark) and T6 reaches $0.768 \pm 0.088$ (higher mean but with one badly underperforming seed). Both graph variants tie V0 within seed variance and remain almost ten points below the fingerprint-only control, so the methodological point above stands: on BACE the graph encoder does not contribute a detectable test-time signal once a strong fingerprint baseline is present.

\subsection{Ablations}
\label{sec:ablations}
We isolate three design choices of V2-T5/T6 that we found most consequential during development.

\noindent\textbf{(a) Bias initialization $b{=}0$ vs.\ $b{=}{+}5$.}
The bias term $b$ in \eqref{eq:v2t5gate} determines the initial value of every edge gate. With $b{=}0$ the gate starts at $\sigma(0){=}0.5$ across all bonds, halving the effective Laplacian weight and pushing the spectrum far from the V0 starting point; we observed in early runs that this configuration produces a noticeably worse contrastive loss in the first $20$ pretraining epochs and never fully recovers to the $b{=}{+}5$ trajectory under the same budget. With $b{=}{+}5$ the gate starts at $\sigma(5)\approx 0.993$, so the epoch-zero filter is, by construction, within rounding of the V0 filter, which stabilizes the first few pretraining epochs. Where the gate ends up after training is a different question. A post-hoc audit of the three BACE V2-T5 checkpoints (\S\ref{sec:experiments}) reveals that the bias term itself moves only a small amount ($b$ ends at $4.77\pm 0.01$ across the three seeds), but the MLP weights learn to add a substantial negative shift to the pre-sigmoid scalar (mean $-2.62 \pm 0.34$ across seeds), so the final per-bond gate distribution is concentrated at $\sigma\approx 0.07$ with seed-level means in $[0.057, 0.101]$. The per-seed gate histograms are unimodal and right-skewed (Figure~\ref{fig:gate}); even the per-seed 99th percentile gate stays under $0.20$ and the maximum stays under $0.27$. In other words, V2-T5 is not learning to route a subset of bonds through (binary on/off) but rather to apply a graded bond-sparsification under which every bond is heavily damped while a long tail of bonds gets several times the typical weight. The near-identity initialization is what we initialize at and not what we end at: the model converges to an aggressively bond-sparsified Laplacian whose effective spectrum is far from V0's. The $b{=}{+}5$ design choice should therefore be read as a stability constraint at epoch zero rather than a guarantee that the trained encoder stays close to V0.

A matched-seed sweep at $b{=}{+}1$ ($\sigma(1)\approx 0.731$) keeps every other hyperparameter identical, produces test AUC $0.759 \pm 0.019$ on BACE---tied with the $b{=}{+}5$ trajectory ($0.763 \pm 0.010$)---and lands on an even more sparsified gate distribution (mean $\sigma\approx 0.04$). So on BACE, bias initialization moves the trained gates but does not move test AUC: fingerprint saturation pins all graph variants to $\sim$$0.76$ regardless of which gate regime the encoder finds. The same ablation on BBBP behaves differently: $b{=}{+}1$ drops test AUC by $3.3$ points (to $0.829 \pm 0.015$) and the per-seed gate distributions tighten to $\{0.096, 0.183, 0.167\}$, none reaching the very-sparse regime that $b{=}{+}5$ produces. Bias initialization therefore matters only when the benchmark is not fingerprint-saturated.

\noindent\textbf{Chemistry-aware routing on BBBP, and only on BBBP.}
Bond-level attribution sharpens this picture. We dump the per-bond gate $g_{ij}$ for every chemical-bond edge in one inference pass over the pretrain split and correlate it with the bond-feature dominant column ($\in$ \{Single, Double, Triple, Aromatic\}). On BACE, $g_{ij}$ is only weakly correlated with bond type at either initialization (Pearson with Aromatic indicator in $[+0.06, +0.18]$ at $b{=}{+}5$ and $[+0.06, +0.17]$ at $b{=}{+}1$; aromatic-vs-non-aromatic mean-gate deltas within noise). On BBBP at $b{=}{+}5$ the same correlation is large and stable: Pearson with Aromatic is $\{+0.690, +0.664, +0.618\}$ across the three seeds (with mirror anti-correlation $\{-0.49, -0.45, -0.39\}$ on Single bonds). The seed-19 BBBP outlier whose gate distribution did not sparsify (mean $\sigma\approx 0.29$, the only BBBP seed where V2-T5 loses V0) shows the same routing \emph{direction}---Aromatic Pearson $0.66$---just at a uniformly higher magnitude; basin heterogeneity is in scale, not direction. Switching the BBBP runs to $b{=}{+}1$ flips this: Pearson with Aromatic collapses to $\{-0.03, -0.02, +0.07\}$ and the model instead up-weights Double bonds (Pearson $\{+0.39, +0.42, +0.44\}$ on $b{=}{+}1$ vs $\{-0.32, -0.34, -0.36\}$ on $b{=}{+}5$). So $b{=}{+}5$ does not just buy basin diversity---it reliably steers the optimizer into the aromatic-preserving basin specifically, which is the chemistry-relevant one for blood-brain barrier permeability (a property driven by through-space $\pi$-system and aromatic-ring lipophilicity), and which delivers the higher test AUC. BACE binding by contrast sits in a single fingerprint-encoded pocket, so the graph routing carries no signal that the Morgan fingerprint isn't already capturing, and neither basin matters at test time.

\noindent\textbf{Sampling variability of the BBBP V2-T5 win.}
The BBBP cells in Table~\ref{tab:main} are now reported at $n{=}9$ training seeds ($\{9, 19, 29, 39, 49, 59, 69, 79, 89\}$); the three deep variants tie within seed variance (V0 $0.828 \pm 0.027$, V2-T5 $0.831 \pm 0.032$, T6 $0.831 \pm 0.022$). The original $n{=}3$ subsample (seeds $9, 19, 29$) placed V2-T5 at $0.862 \pm 0.025$, which we now read as upper-tail sampling: the per-seed means across all nine seeds are $\{0.862, 0.838, 0.887, 0.798, 0.817, 0.815, 0.781, 0.838, 0.841\}$, and the original three rank in the top $3/9, 5/9, 1/9$. V0 and T6 extensions move in the opposite direction (V0 drops from $0.836$ to $0.828$; T6 drops from $0.840$ to $0.831$), and the matched-seed sign test on the original three seeds no longer reflects a population-level effect: at $n{=}9$ the V2-T5-over-V0 mean delta is only $+0.003$. The chemistry-aware aromatic-routing finding from the previous paragraph is robust across all six audited V2-T5 seeds (the routing direction is identical and the Pearson correlation between gate and aromatic-bond indicator stays in $[0.6, 0.7]$), so the more reliable claim from this paper is the mechanism rather than the headline AUC number.

\begin{figure}[t]
\centering
\includegraphics[width=\columnwidth]{gate_distribution_bace_v2t5.pdf}
\caption{Trained per-bond gate distribution for V2-T5 on BACE, across three pretraining seeds and two bias-initialization choices. Blue group: $b{=}{+}5$ initialization ($\sigma(5)\approx 0.993$, the default in the paper). Red group: $b{=}{+}1$ ablation ($\sigma(1)\approx 0.731$). Dotted lines mark the initial gate values for each group. In both groups the bias term barely moves during training (final $b\approx 4.77$ for the $+5$ runs and $b\approx 0.83$ for the $+1$ runs), but the MLP weights learn substantial negative shifts, so the trained gates sit far below the initialization. The $b{=}{+}1$ trajectory converges to an even more sparsified Laplacian (mean $\sigma\approx 0.04$); on BACE it ties the $b{=}{+}5$ default at test time ($0.759 \pm 0.019$ vs.\ $0.763 \pm 0.010$, fingerprint saturated), while on BBBP it drops $3.3$ AUC points and the gate fails to enter the aromatic-routing basin.}
\label{fig:gate}
\end{figure}

\noindent\textbf{(b) Uni-Mol pair vs.\ random pair representation.}
The static gate of V2-T5 receives a frozen Uni-Mol pair tensor as its only three-dimensional signal. A natural concern is whether the gain over V0 reflects geometric information specifically, or merely the additional parameters introduced by the gating MLP. We address this by replacing $\mathbf{P}$ with an isotropic Gaussian tensor of identical shape (the gate retains its $b{=}{+}5$ bias, the MLP retains its capacity) and re-running the BBBP pipeline on the original three matched seeds. The random-pair variant reaches a per-seed-mean AUC of $0.848 / 0.814 / 0.828$ on seeds $9 / 19 / 29$ ($0.830 \pm 0.017$ overall), losing to the Uni-Mol-pair V2-T5 on all three matched seeds (deltas $+0.014, +0.024, +0.059$; mean delta $+0.032$, $3/3$ wins, one-sided binomial $p{=}0.125$). So at $n{=}3$ the Uni-Mol pair tensor does carry a small per-seed signal that an isotropic random tensor of the same shape does not capture, even though---per Table~\ref{tab:main} and \S\ref{sec:experiments}---neither variant moves above the V0 envelope once the comparison is taken to $n{=}9$. We read this as: on BBBP, the V2-T5 architecture functions primarily as a structured MLP gate; the Uni-Mol geometry contributes a measurable per-seed effect that is below the seed-to-seed test-AUC noise floor of the benchmark.

\noindent\textbf{(c) Static gate (V2-T5) vs.\ dynamic update (T6).}
At initialization, T6 is strictly more expressive than V2-T5 because its zero-initialized pair-update block reduces to the identity, leaving the static gate as a special case. In practice, however, T6 underperforms V2-T5 on BBBP by $2.2$ AUC points on average and shows higher seed-level variance ($0.036$ vs.\ $0.025$). The variance is concentrated in seed $19$, where T6 produces a mean AUC of $0.800$ while V2-T5 produces $0.838$ on the identical split; we interpret this as evidence that the additional optimization degrees of freedom in T6's pair-update block, in the small-data regime of BBBP, occasionally fail to find an update direction that the contrastive objective rewards consistently across the four views. We return to this point in Section~\ref{sec:discussion}.

% ---------------------------------------------------------------------
\section{Discussion}
\label{sec:discussion}

\noindent\textbf{Why static beats dynamic on small benchmarks.}
A natural expectation, given the relative parameter counts, is that the more expressive T6 should dominate V2-T5 on every benchmark, with the gap widening on the smallest datasets where 3D information is most needed. The opposite happens. We attribute this inversion to two effects that compound in the small-data regime. First, the contrastive pretraining objective is not directly aligned with the downstream property: a richer update mechanism gives the encoder more ways to satisfy NT-Xent that do not translate into a useful feature for AUC. Second, with $\sim$2000 molecules and a single training-validation split, the seed-to-seed variance of the linear-probe estimate is large enough ($\pm 0.025$--$0.036$ AUC) that a model with even slightly higher optimization variance can be reliably ranked below a more conservative model. V2-T5's near-identity initialization is a strong form of starting-point conservatism, even though, as the BACE audit in \S\ref{sec:v2t5} shows, the trained gate converges to $\sigma\approx 0.07$ rather than to the $\sigma\approx 0.993$ where it started. We read this as: the $b{=}{+}5$ initialization buys a stable first-few-epoch regime under contrastive pretraining, and from there the model is free to slide into a sparser Laplacian. The reason this still produces a more conservative encoder than T6 is that the per-bond scalar gate of V2-T5 is a much lower-capacity perturbation than T6's bidirectional pair-node attention update: T6 can re-shape the pair tensor itself at every layer, whereas V2-T5 can only multiply each bond by a single sigmoid. T6's degradation on BBBP is also concentrated in a single unlucky seed rather than uniformly distributed, which is consistent with higher per-seed optimization variance from that additional capacity. A related symptom in the V0 baseline supports the variance argument: the validation-best epoch on the three BBBP V0 seeds is $233$, $767$, and $332$, with the seed that picks epoch $767$ also being the one on which V2-T5 loses to V0 (\S\ref{sec:experiments}). Late-checkpoint selection on small benchmarks is itself a non-trivial source of seed-level variance, independent of any encoder difference.

\noindent\textbf{What the fingerprint-only control tells us.}
The fingerprint-only control achieves $0.846$ AUC on BACE while the V0 graph encoder achieves $0.757$, and the gap is large relative to the seed standard deviation of either ($0.012$ vs.\ $0.047$). A still stronger classical baseline---a 300-tree random forest on Morgan fingerprints ($1024$ bits, radius $2$) concatenated with eleven \texttt{RDKit} two-dimensional descriptors, matched seed-for-seed and split-for-split to the deep models---reaches $0.895 \pm 0.003$ AUC on the same Uni-Mol fold split when averaged over the 30-seed protocol of Deng et al.~\cite{deng2023systematic}, strictly exceeding the deep fingerprint MLP and every deep graph variant on BACE (V2-T5 $0.763 \pm 0.010$, T6 $0.768 \pm 0.088$) by margins many seed-standard-deviations wide. We read this not as a negative result for graph neural networks but as a methodological flag, and one that is now also documented at scale: Deng et al.~\cite{deng2023systematic} train $62{,}820$ model variants and report that the random forest statistically ties or beats the strongest pretrained GNNs on most of the \texttt{MoleculeNet} classification suite, with BACE specifically among the datasets where the simple classical model wins. A fingerprint MLP is cheap, has been a strong baseline since the late 2000s, and is often missing from molecular property prediction tables that compare expensive graph architectures against each other; the same is true a fortiori for the random forest. For completeness we also report Chemprop~\cite{yang2019chemprop}, the standard supervised D-MPNN baseline, under the same matched protocol; on BBBP and FreeSolv it ranks below V2-T5 (Table~\ref{tab:main}), and on BACE it sits well above V2-T5 ($0.840$ vs $0.763$) but still well below the random forest ($0.895$), consistent with the same fingerprint-saturation reading rather than a graph-architecture deficit. On BBBP, V2-T5 at $n{=}9$ ($0.831 \pm 0.032$) sits a single seed-standard-deviation above the random-forest control ($0.799 \pm 0.007$) and above Chemprop ($0.815 \pm 0.014$), while tying both other deep variants---a more modest test-AUC story than the $n{=}3$ subsample suggested, but, in our view, only credible at all because it is measured against a protocol in which all three reference families---classical fingerprint, classical fingerprint-plus-trees, and supervised graph---are reported alongside the deep models.

\noindent\textbf{Limitations.}
We evaluate on three small benchmarks. The largest dataset for which we report numbers, BBBP, contains $2{,}039$ molecules, and so all claims should be read as evidence about small-to-moderate molecular benchmarks rather than about large-scale screening; Wijaya et al.~\cite{wijaya2024twostage} report that on small molecular datasets pretrained models often show a significant drop in performance, sometimes only marginally surpassing a baseline that simply predicts the mean label value, and our random-forest control on BACE is consistent with that observation. The $n{=}3$-seed reporting protocol we inherit from the Uni-Mol scaffold-fold pipeline is itself a fragility: the dedicated $n{=}9$ extension we report in \S\ref{sec:experiments} shows that V2-T5's $+2.6$ AUC margin on BBBP at $n{=}3$ does not survive replication, even though the underlying chemistry-aware aromatic-routing mechanism does (every audited seed up-weights aromatic bonds). Readers should treat the test-AUC win as $n{=}3$-sized evidence and the gate-attribution finding as the more reliable claim from this paper. We also rely on a scaffold-style split which, as Guo et al.~\cite{guo2024scaffold} document on 60 NCI-60 cancer-cell-line datasets, tends to produce structurally similar train/test pairs and therefore overestimate downstream-task utility relative to clustering- or similarity-aware splits. The community has been migrating to benchmarks specifically designed to address these shortcomings---\texttt{Therapeutics Data Commons}~\cite{huang2021tdc} for ADMET endpoints and the \texttt{Polaris} hub~\cite{tossou2024polaris} for industry-curated OOD splits---and our claims should be re-tested under those protocols before being generalized; a recent reproducibility audit~\cite{tdcaudit2026} additionally cautions that several top TDC leaderboard entries fail reproducibility checks, so even the migration target requires care. The three-dimensional signal comes from a single pretrained encoder (Uni-Mol); we did not test whether an alternative 3D encoder (DimeNet, EGNN) plugged into the same V2-T5 gate would produce comparable gains. The HIV benchmark ($\sim$41{,}000 molecules) is the natural next step and is in onboarding but not yet reported; similarly, multi-label tasks such as ClinTox and Tox21 are deferred to a later phase. Finally, all BACE graph-variant numbers we report come from a corrected data rebuild that fixed a directed-edge bug in the bond featurization; we deliberately do not report the numbers from the earlier, defective pipeline, because that data-graph defect would violate the fair-comparison protocol the rest of the paper relies on.

% ---------------------------------------------------------------------
\section{Conclusion}
\label{sec:conclusion}

We have shown that a frozen three-dimensional pair representation can be injected into a Chebyshev~II spectral GNN through a single scalar gate per chemical bond, and that the resulting V2-T5 gate, even when initialized near identity ($\sigma(b{=}{+}5)\approx 0.993$), converges during training to an aggressively bond-sparsified distribution. On BBBP that gate is chemistry-aware: across every audited seed, V2-T5 up-weights aromatic bonds and down-weights saturated single bonds, with Pearson correlations to the aromatic-bond indicator in the $[0.6, 0.7]$ band. On BACE the same architecture, same initialization, and same contrastive objective learn no such routing, and the test AUC of every V2-T5 seed sits at the $\sim$$0.76$ ceiling that the random forest already exceeds by $13$ points. The win in test AUC is therefore benchmark-dependent and, at $n{=}3$ training seeds, within sampling noise on BBBP; the chemistry-aware routing mechanism, by contrast, is the more reliable empirical claim, holds across all six audited V2-T5 checkpoints, and offers a concrete interpretability handle for spectral GNN gating in molecular property prediction. A fingerprint-only control turns out to be the most competitive single-modality baseline on BACE---evidence, in our reading, that the molecular property prediction literature should report such controls more routinely. Three directions stand out for follow-up work. First, scaling the same gate-with-near-identity recipe to HIV ($\sim$41{,}000 molecules) and to multi-label tasks (ClinTox, Tox21) will test whether the static design continues to dominate once the small-data variance argument no longer applies. Second, allowing the upstream 3D encoder to receive gradient signal from the downstream loss---rather than treating Uni-Mol as a frozen feature extractor---would let the pair representation specialize to the spectral filter's actual needs. Third, the low-rank tensor-gate generalization sketched at the end of \S\ref{sec:v2t5}---one scalar per Chebyshev tower per bond rather than one scalar per bond---would let the same architecture modulate the high-pass and low-pass paths independently, which we believe is a promising next step.

% ---------------------------------------------------------------------
\begin{thebibliography}{00}
\bibitem{wang2022molclr} Y. Wang, J. Wang, Z. Cao, and A. Barati Farimani, ``Molecular contrastive learning of representations via graph neural networks,'' \textit{Nat. Mach. Intell.}, vol. 4, no. 3, pp. 279--287, 2022.
\bibitem{rong2020grover} Y. Rong et al., ``Self-supervised graph transformer on large-scale molecular data,'' in \textit{Proc. Adv. Neural Inf. Process. Syst. (NeurIPS)}, 2020.
\bibitem{zhou2023unimol} G. Zhou et al., ``Uni-Mol: A universal 3D molecular representation learning framework,'' in \textit{Proc. Int. Conf. Learn. Represent. (ICLR)}, 2023.
\bibitem{liu2022graphmvp} S. Liu et al., ``Pre-training molecular graph representation with 3D geometry,'' in \textit{Proc. Int. Conf. Learn. Represent. (ICLR)}, 2022.
\bibitem{he2022chebnetii} M. He, Z. Wei, and J.-R. Wen, ``Convolutional neural networks on graphs with Chebyshev approximation, revisited,'' in \textit{Proc. Adv. Neural Inf. Process. Syst. (NeurIPS)}, 2022.
\bibitem{defferrard2016chebnet} M. Defferrard, X. Bresson, and P. Vandergheynst, ``Convolutional neural networks on graphs with fast localized spectral filtering,'' in \textit{Proc. Adv. Neural Inf. Process. Syst. (NeurIPS)}, 2016.
\bibitem{kipf2017gcn} T. N. Kipf and M. Welling, ``Semi-supervised classification with graph convolutional networks,'' in \textit{Proc. Int. Conf. Learn. Represent. (ICLR)}, 2017.
\bibitem{velickovic2018gat} P. Veli\v{c}kovi\'{c} et al., ``Graph attention networks,'' in \textit{Proc. Int. Conf. Learn. Represent. (ICLR)}, 2018.
\bibitem{xu2019gin} K. Xu, W. Hu, J. Leskovec, and S. Jegelka, ``How powerful are graph neural networks?'' in \textit{Proc. Int. Conf. Learn. Represent. (ICLR)}, 2019.
\bibitem{he2021bernnet} M. He, Z. Wei, H. Xu, and J.-R. Wen, ``BernNet: Learning arbitrary graph spectral filters via Bernstein approximation,'' in \textit{Proc. Adv. Neural Inf. Process. Syst. (NeurIPS)}, 2021.
\bibitem{chien2021gprgnn} E. Chien, J. Peng, P. Li, and O. Milenkovic, ``Adaptive universal generalized PageRank graph neural network,'' in \textit{Proc. Int. Conf. Learn. Represent. (ICLR)}, 2021.
\bibitem{shang2021eagcn} C. Shang et al., ``Multi-view spectral graph convolution with consistent edge attention for molecular modeling,'' \textit{Neurocomputing}, vol. 445, pp. 12--25, 2021.
\bibitem{chen2024polygcl} J. Chen, R. Lei, and Z. Wei, ``PolyGCL: Graph contrastive learning via learnable spectral polynomial filters,'' in \textit{Proc. Int. Conf. Learn. Represent. (ICLR)}, 2024.
\bibitem{nogueira2025spectra} B. Nogueira, M. Jiang, N. V. Chawla, and N. Moniz, ``SPECTRA: Spectral target-aware graph augmentation for imbalanced molecular property regression,'' arXiv:2511.04838, 2025.
\bibitem{gilmer2017mpnn} J. Gilmer, S. S. Schoenholz, P. F. Riley, O. Vinyals, and G. E. Dahl, ``Neural message passing for quantum chemistry,'' in \textit{Proc. Int. Conf. Mach. Learn. (ICML)}, 2017.
\bibitem{xiong2020attentivefp} Z. Xiong et al., ``Pushing the boundaries of molecular representation for drug discovery with the graph attention mechanism,'' \textit{J. Med. Chem.}, vol. 63, no. 16, pp. 8749--8760, 2020.
\bibitem{schutt2017schnet} K. T. Sch\"{u}tt et al., ``SchNet: A continuous-filter convolutional neural network for modeling quantum interactions,'' in \textit{Proc. Adv. Neural Inf. Process. Syst. (NeurIPS)}, 2017.
\bibitem{klicpera2020dimenet} J. Klicpera, J. Gro\ss{}, and S. G\"{u}nnemann, ``Directional message passing for molecular graphs,'' in \textit{Proc. Int. Conf. Learn. Represent. (ICLR)}, 2020.
\bibitem{satorras2021egnn} V. G. Satorras, E. Hoogeboom, and M. Welling, ``E(n) equivariant graph neural networks,'' in \textit{Proc. Int. Conf. Mach. Learn. (ICML)}, 2021.
\bibitem{feng2024frad} S. Feng, Y. Ni, Y. Lan, Z.-M. Ma, and W.-Y. Ma, ``Fractional denoising for 3D molecular pre-training,'' \textit{Nat. Mach. Intell.}, 2024.
\bibitem{fang2022gem} X. Fang et al., ``Geometry-enhanced molecular representation learning for property prediction,'' \textit{Nat. Mach. Intell.}, vol. 4, no. 2, pp. 127--134, 2022.
\bibitem{stark2022infomax} H. St\"{a}rk et al., ``3D Infomax improves GNNs for molecular property prediction,'' in \textit{Proc. Int. Conf. Mach. Learn. (ICML)}, 2022.
\bibitem{xia2023molebert} J. Xia, Y. Zhu, Y. Du, Y. Liu, and S. Z. Li, ``Mole-BERT: Rethinking pre-training graph neural networks for molecules,'' in \textit{Proc. Int. Conf. Learn. Represent. (ICLR)}, 2023.
\bibitem{li2023kpgt} H. Li, D. Zhao, and J. Zeng, ``A knowledge-guided pre-training framework for improving molecular representation learning,'' \textit{Nat. Commun.}, vol. 14, art. no. 7568, 2023.
\bibitem{klaeser2024minimol} K. Kl\"{a}ser et al., ``MiniMol: A parameter-efficient foundation model for molecular learning,'' arXiv:2404.14986, 2024.
\bibitem{notwell2023maplight} J. H. Notwell and M. W. Wood, ``ADMET property prediction through combinations of molecular fingerprints,'' arXiv:2310.00174, 2023.
\bibitem{novoexpert2026} NovoExpert Team, ``NovoExpert-2: State-of-the-art ADMET prediction via gradient-boosted trees on MapLight fingerprints and GIN embeddings,'' ChemRxiv preprint, doi: 10.26434/chemrxiv.15000061/v2, 2026.
\bibitem{wu2018moleculenet} Z. Wu et al., ``MoleculeNet: A benchmark for molecular machine learning,'' \textit{Chem. Sci.}, vol. 9, no. 2, pp. 513--530, 2018.
\bibitem{deng2023systematic} J. Deng et al., ``A systematic study of key elements underlying molecular property prediction,'' \textit{Nat. Commun.}, vol. 14, art. no. 6395, 2023.
\bibitem{yang2019chemprop} K. Yang et al., ``Analyzing learned molecular representations for property prediction,'' \textit{J. Chem. Inf. Model.}, vol. 59, no. 8, pp. 3370--3388, 2019.
\bibitem{wijaya2024twostage} K. Wijaya et al., ``Two-stage pretraining for molecular property prediction in the wild,'' in \textit{Proc. ICML AI for Science Workshop}, 2024, arXiv:2411.03537.
\bibitem{guo2024scaffold} Q. Guo, S. Hernandez-Hernandez, and P. J. Ballester, ``Scaffold splits overestimate virtual screening performance,'' in \textit{Proc. Int. Conf. Artif. Neural Netw. (ICANN)}, ser. LNCS 15025, 2024, pp. 58--72.
\bibitem{huang2021tdc} K. Huang et al., ``Therapeutics Data Commons: Machine learning datasets and tasks for drug discovery and development,'' in \textit{Proc. NeurIPS Datasets and Benchmarks Track}, 2021.
\bibitem{tossou2024polaris} P. Tossou, C. Wognum, M. Craig, H. Mary, and E. Noutahi, ``Real-world molecular out-of-distribution: Specification and investigation,'' \textit{J. Chem. Inf. Model.}, 2024.
\bibitem{tdcaudit2026} TDC Audit Authors, ``Critical assessment of ML models for ADMET prediction in TDC leaderboards,'' bioRxiv preprint, doi: 10.64898/2026.02.26.708193, 2026.
\end{thebibliography}

\end{document}
